\documentclass[twocolumn,superscriptaddress,unsortedaddress,nofootinbib,amsmath,amssymb,
aps,floatfix]{revtex4-2}
\usepackage{graphicx,tikz} 
\usetikzlibrary{arrows.meta,decorations.pathreplacing,calc,positioning}
\usepackage[mode=buildmissing]{standalone} 
\usepackage{dcolumn}
\usepackage{bm}
\usepackage{physics}
\usepackage{xcolor}
\usepackage{siunitx}
\usepackage{soul}
\usepackage{comment}
\usepackage[bb=boondox]{mathalfa}
\usetikzlibrary{external}
\tikzexternalize[prefix=figures/]
\DeclareMathOperator{\sinc}{sinc}
\newcommand{\may}[1]{\textcolor{orange}{#1}}
\begin{document}
\title{Lipkin model, Entanglement, Noise, Dephasing and dissipation}

\author{Juan David Álvarez-Cuartas}
\affiliation{Departament of Physics and Center for Bioinformatics and Photonics---CIBioFi, 
Universidad del Valle, Cali 760032, Colombia}
\affiliation{Department of Physics and Center for Computing in Science Education, University of Oslo, N-0316 Oslo, Norway}
\author{Patrick Cook and Danny Jammooa}
\affiliation{Facility for Rare Isotope Beams and Department of Physics and Astronomy, Michigan State University, East Lansing, Michigan 48824, USA}
\author{Morten Hjorth-Jensen}
\affiliation{Department of Physics and Center for Computing in Science Education, University of Oslo, N-0316 Oslo, Norway}
\author{Dean Lee}
\affiliation{Facility for Rare Isotope Beams and Department of Physics and Astronomy, Michigan State University, East Lansing, Michigan 48824, USA}
\author{John H.~Reina, and Juan M.~Scarpetta}
\affiliation{Departament of Physics and Center for Bioinformatics and Photonics---CIBioFi, 
Universidad del Valle, Cali 760032, Colombia}
\affiliation{Department of Physics and Center for Computing in Science Education, University of Oslo, N-0316 Oslo, Norway}
\begin{abstract}
Text to come here
\end{abstract}

%\pacs{02.70.Ss, 31.15.A-, 31.15.bw, 71.15.-m, 73.21.La}


\date{\today}

\maketitle

\section{Introduction} 
Structure of article
\begin{enumerate}
    \item Introduction with motivation
    \item Discussion of the Lipkin model
    \item Setting up time evolution
    \item PMMs and time evolution
    \item 
\end{enumerate}

\section{Mathematical Framework and Entanglement Measures}

\subsection{Hamiltonian and Symmetry.}
The Lipkin--Meshkov--Glick (LMG) model describes $N$ spin-$\tfrac{1}{2}$ particles interacting through collective spin operators.  
For four fermions, the case we consider in the examples  here, each levels has
degeneration $d=4$, leading to different total spin values.  The two
levels have quantum numbers $\sigma=\pm 1$, with the upper level
having $2\sigma=+1$ and energy $\varepsilon_{1}= \varepsilon/2$. The
lower level has $2\sigma=-1$ and energy
$\varepsilon_{2}=-\varepsilon/2$. That is, the lowest single-particle
level has negative spin projection (or spin down), while the upper
level has spin up.  In addition, the substates of each level are
characterized by the quantum numbers $p=1,2,3,4$.



The Hamiltonian of the system is given by
\[
\begin{array}{ll}
\hat{H}=&\hat{H}_{0}+\hat{H}_{1}+\hat{H}_{2}\\
&\\
\hat{H}_{0}=&\frac{1}{2}\varepsilon\sum_{\sigma ,p}\sigma
a_{\sigma,p}^{\dagger}a_{\sigma ,p}\\
&\\
\hat{H}_{1}=&\frac{1}{2}V\sum_{\sigma ,p,p'}
a_{\sigma,p}^{\dagger}a_{\sigma ,p'}^{\dagger}
a_{-\sigma ,p'}a_{-\sigma ,p}\\
&\\
\hat{H}_{2}=&\frac{1}{2}W\sum_{\sigma ,p,p'}
a_{\sigma,p}^{\dagger}a_{-\sigma ,p'}^{\dagger}
a_{\sigma ,p'}a_{-\sigma ,p}\\
&\\
\end{array}
\]
where $V$ and $W$ are constants. The operator 
$H_{1}$ can move pairs of fermions
while $H_{2}$ is a spin-exchange term. The latter
moves a pair of fermions from a state $(p\sigma ,p' -\sigma)$ to a state
$(p-\sigma ,p'\sigma)$.

We are going to rewrite the above Hamiltonian in terms of so-called  quasispin operators
\[
\begin{array}{ll}
\hat{J}_{+}=&\sum_{p}
a_{p+}^{\dagger}a_{p-}\\
&\\
\hat{J}_{-}=&\sum_{p}
a_{p-}^{\dagger}a_{p+}\\
&\\
\hat{J}_{z}=&\frac{1}{2}\sum_{p\sigma}\sigma
a_{p\sigma}^{\dagger}a_{p\sigma}\\
&\\
\hat{J}^{2}=&J_{+}J_{-}+J_{z}^{2}-J_{z}\\
&\\
\end{array}
\]
and the number operator 
\[
\hat{N}=\sum_{p\sigma}
a_{p\sigma}^{\dagger}a_{p\sigma}.
\]

We will now rewrite the Hamiltonian in terms of the above quasi-spin operators and the number operator 
\begin{equation}
N = \sum_{p,\sigma} a_{p\sigma}^\dagger a_{p\sigma}.
label{eq:N}
\end{equation}
Going through each term of the Hamiltonian and using the expressions for the quasi-spin operators we obtain
\begin{equation}
H_0 = \varepsilon J_z, 
\end{equation}
\begin{equation}
H_1 = \frac{1}{2} V \left( J_+^2 + J_-^2 \right),
\end{equation}
and 
\begin{equation}
H_2 = \frac{1}{2} W \left( -N + J_+ J_- + J_- J_+ \right).
\end{equation}

The above expressions can in turn be used to show that the Hamiltonian
commutes with the various quasi-spin operators. This leads to quantum
numbers which are conserved.  Let us first show that $[H,J^2]=0$,
which means that $J$ is a so-called *good* quantum number and that the
total spin is a conserved quantum number.

The matrix elements are given by $\langle J,J_z \vert H \vert J',J_z' \rangle$.
The
Hamiltonian is hermitian and we obtain after all this labor of ours

\begin{equation}
H_{J = 2} =
\begin{bmatrix}
-2\varepsilon & 0 & \sqrt{6}V & 0 & 0 \\
0 & -\varepsilon + 3W & 0 & 3V & 0 \\
\sqrt{6}V & 0 & 4W & 0 & \sqrt{6}V \\
0 & 3V & 0 & \varepsilon + 3W & 0 \\
0 & 0 & \sqrt{6}V & 0 & 2\varepsilon
\end{bmatrix}
label{eq:HJ=2}
\end{equation}





Its Hamiltonian can be expressed as
\begin{equation}
  \hat{H}_{\mathrm{LMG}} = -\frac{\lambda}{N} \left( \hat{S}_x^2 + \gamma\, \hat{S}_y^2 \right)
  - h\, \hat{S}_z,
  \label{eq:lmg-hamiltonian}
\end{equation}
where $\hat{S}_\alpha = \tfrac{1}{2}\sum_{i=1}^N \hat{\sigma}_i^\alpha$ are collective spin operators,
$\lambda$ is the interaction strength, $\gamma$ the anisotropy parameter, and $h$ an external magnetic field.
The factor $1/N$ ensures a well-defined thermodynamic limit.
The total spin $\hat{\mathbf{S}}^2$ commutes with the Hamiltonian, implying conservation of total spin $S = N/2$,
and the system therefore lives in the symmetric subspace spanned by the Dicke states
\begin{equation}
  \ket{S, M} \equiv \ket{N/2, M}, \quad M = -S, \ldots, S.
\end{equation}




\subsection{Mean-Field and Quantum Phase Transition.}
In the thermodynamic limit ($N \to \infty$), a mean-field treatment becomes exact.
A semiclassical energy functional can be written in terms of collective spin variables
$(\theta, \phi)$ on the Bloch sphere as
\begin{equation}
  E(\theta, \phi) = -\frac{\lambda}{2}\sin^2\theta \left( \cos^2\phi + \gamma\, \sin^2\phi \right)
  - h\, \cos\theta.
\end{equation}
The ground state undergoes a second-order quantum phase transition at $h_c = \lambda(1+\gamma)/2$,
separating a symmetric phase ($\langle S_x \rangle = 0$) from a symmetry-broken phase ($\langle S_x \rangle \neq 0$).
Near this critical point, the entanglement entropy exhibits logarithmic scaling with system size, signalling a diverging correlation length in the collective degrees of freedom.

\subsection{Reduced Density Matrices and Entropy.}
Because of full permutation symmetry, the reduced density matrix of any single spin
is identical and can be obtained from the collective spin density operator $\rho$ as
\begin{equation}
  \rho_1 = \mathrm{Tr}_{2,\ldots,N}(\rho)
  = \frac{1}{2}\left( \mathbb{I} + \frac{2}{N}\langle \hat{\mathbf{S}} \rangle \cdot \boldsymbol{\sigma} \right).
\end{equation}
The single-spin entanglement (or \emph{one-tangle}) with the remainder of the system is quantified by the linear entropy
\begin{equation}
  \tau_1 = 4 \det(\rho_1) = 1 - \frac{4}{N^2} \langle \hat{\mathbf{S}} \rangle^2.
\end{equation}
This quantity vanishes in product states and reaches unity for maximally mixed single-spin states.

For bipartite entanglement between two subsystems of sizes $L$ and $N-L$, the von Neumann entropy
\begin{equation}
  S_L = - \mathrm{Tr}\, (\rho_L \log_2 \rho_L)
\end{equation}
is the standard measure.  In the symmetric Dicke basis, $\rho_L$ is block-diagonal and
its eigenvalues can be expressed analytically in terms of Clebsch--Gordan coefficients,
allowing explicit evaluation of $S_L$ for moderate $N$.

\subsection{Two-Spin Correlations and Concurrence.}
Pairwise entanglement between two spins is characterized by the \emph{concurrence} $C$,
defined from the two-spin reduced density matrix $\rho_{12}$ as
\begin{equation}
  C = \max \left\{ 0, \lambda_1 - \lambda_2 - \lambda_3 - \lambda_4 \right\},
\end{equation}
where $\lambda_i$ are the square roots of the eigenvalues (in descending order)
of the matrix $\rho_{12} \tilde{\rho}_{12}$, with
$\tilde{\rho}_{12} = (\sigma_y \otimes \sigma_y)\, \rho_{12}^*\, (\sigma_y \otimes \sigma_y)$.
In the LMG model, due to full permutation symmetry, $C$ is identical for all spin pairs.
It reaches a maximum near the critical point and decays as $N^{-1/3}$ at criticality,
indicating that bipartite entanglement is sub-extensive even though the total
multipartite entanglement diverges.

\subsection{Multipartite and Mixed-State Measures.}
Beyond bipartite partitions, the \emph{global entanglement} and the
\emph{geometric measure of entanglement} are often used:
\begin{equation}
  E_{\mathrm{G}} = 1 - \max_{\ket{\phi_{\mathrm{sep}}}} |\langle \phi_{\mathrm{sep}} | \psi_0 \rangle|^2,
\end{equation}
where $\ket{\phi_{\mathrm{sep}}}$ is the closest separable state to the ground state $\ket{\psi_0}$.
For mixed or non-complementary partitions, the logarithmic negativity
\begin{equation}
  \mathcal{N} = \log_2 \| \rho^{T_A} \|_1
\end{equation}
quantifies entanglement by means of the trace norm of the partially transposed density matrix $\rho^{T_A}$.

\subsection{Scaling and Universality.}
Near the quantum critical point, different entanglement measures show universal scaling with system size $N$:
\begin{equation}
  S_L \sim \frac{1}{6} \log_2 N + \mathrm{const}, \qquad
  \tau_1 \sim N^{-2/3}, \qquad
  C \sim N^{-1/3}.
\end{equation}
These exponents coincide with those predicted by conformal field theory
for one-dimensional critical systems, even though the LMG model is infinite-range.
Such universality underscores that entanglement, rather than spatial locality,
captures the essential signatures of quantum criticality.

\section{Results and discussions}


\section{Conclusion}

%\bibliography{paper}
\end{document}


